\chapter{Glossaire}
\label{ann:glossaire}

Ce glossaire reprend sous une forme abrégée les définitions des chapitres. Le renvoi indique le chapitre où le terme est défini et employé.

\begin{description}
    \item[Alignement stratégique] Accord entre la stratégie de l'organisation, celle de son système d'information et leurs organisations respectives (Ch.~\ref{chap:alignement}).
    \item[Analyse des écarts] Comparaison, capacité par capacité, de l'existant et de la cible (Ch.~\ref{chap:cible}).
    \item[Autorisation d'engagement] Limite supérieure des dépenses pouvant être engagées, y compris sur plusieurs années (Ch.~\ref{chap:droit}).
    \item[Bénéfice] Avantage mesurable ou observable obtenu par le métier grâce à un projet. Il a un propriétaire et une mesure de départ (Ch.~\ref{chap:piloter}).
    \item[Bilan d'impact sur l'activité] Étude qui fixe, pour chaque activité, l'interruption et la perte de données admissibles (Ch.~\ref{chap:maitriser}).
    \item[Cadrage] Phase qui répond par écrit aux six questions préalables et se clôt par une note signée (Ch.~\ref{chap:cadrer}).
    \item[Capacité] Aptitude attendue du système d'information, formulée par un verbe et un résultat, sans nommer de solution (Ch.~\ref{chap:cible}).
    \item[Carte d'applicabilité] Tableau daté qui dit, texte par texte, s'il s'applique à l'organisation et sous quelle condition (Ch.~\ref{chap:droit}).
    \item[Cellule de portefeuille] Petite équipe indépendante des projets qui tient les fiches, mesure l'avancement et alerte (Ch.~\ref{chap:piloter}).
    \item[Cible] Description du système d'information voulu à la fin de la période (Ch.~\ref{chap:cible}).
    \item[Comité de pilotage] Instance qui décide pendant la démarche, présidée par celui qui arbitre les budgets (Ch.~\ref{chap:cadrer}).
    \item[Comité de portefeuille] Instance qui lance, arrête et arbitre les projets pendant l'exécution (Ch.~\ref{chap:piloter}).
    \item[Constat] Énoncé en trois temps~: le fait, la pièce, la conséquence pour le métier (Ch.~\ref{chap:diagnostic}).
    \item[Coût complet de possession] Somme de l'investissement, du fonctionnement, du renouvellement et de la sortie sur une période (Ch.~\ref{chap:financer}).
    \item[Crédit de paiement] Limite supérieure des dépenses pouvant être payées dans l'année (Ch.~\ref{chap:droit}).
    \item[Donnée de référence] Donnée partagée par plusieurs processus, dotée d'un propriétaire et d'un système maître (Ch.~\ref{chap:donnees}).
    \item[Durée maximale d'interruption admissible] Délai au-delà duquel l'arrêt d'une activité cause un dommage inacceptable (Ch.~\ref{chap:maitriser}).
    \item[Exigence] Énoncé vérifiable par une pièce, tiré d'un texte, d'une norme, d'un engagement ou d'une règle de l'art (Ch.~\ref{chap:audit}).
    \item[Facteur clé de succès] Domaine où les résultats doivent être bons pour que l'organisation réussisse (Ch.~\ref{chap:alignement}).
    \item[Fiche d'application] Description d'une application en dix rubriques (Ch.~\ref{chap:diagnostic}).
    \item[Fiche de projet] Description d'un projet candidat en une page et huit rubriques (Ch.~\ref{chap:portefeuille}).
    \item[Gouvernance du système d'information] Système par lequel l'usage du numérique est dirigé et contrôlé. Trois tâches~: évaluer, diriger, surveiller (Ch.~\ref{chap:normes}).
    \item[Indice d'alignement prouvé] Part des exigences d'alignement dont la satisfaction est prouvée par une pièce (Ch.~\ref{chap:audit}).
    \item[Informatique hors contrôle] Outils employés sans que la fonction informatique les ait fournis ni recensés (Ch.~\ref{chap:diagnostic}).
    \item[Interopérabilité] Aptitude de systèmes distincts à échanger des données et à les employer (Ch.~\ref{chap:donnees}).
    \item[Lot] Partie d'un grand projet qui livre un résultat utilisable seul (Ch.~\ref{chap:portefeuille}).
    \item[NEK] Cote donnée quand l'auditeur ne peut pas conclure. De preuve si la pièce manque, de critère si l'exigence ne peut pas s'appliquer (Ch.~\ref{chap:audit}).
    \item[Note de cadrage] Document de dix pages, signé du commanditaire, qui fixe la démarche (Ch.~\ref{chap:cadrer}).
    \item[Perte de données maximale admissible] Ancienneté maximale des données que l'on accepte de retrouver après un sinistre (Ch.~\ref{chap:maitriser}).
    \item[Plan fonctionnel] Découpage du système en zones, quartiers et blocs, indépendant des logiciels (Ch.~\ref{chap:cible}).
    \item[Portefeuille] Ensemble des projets réunis pour atteindre des objectifs stratégiques (Ch.~\ref{chap:portefeuille}).
    \item[Principe directeur] Règle durable qui tranche d'avance une famille de décisions. Elle vaut par ses conséquences (Ch.~\ref{chap:cible}).
    \item[Projet obligatoire] Projet imposé par un texte, une injonction ou un risque que l'organisation ne peut pas porter (Ch.~\ref{chap:portefeuille}).
    \item[Provision pour aléas] Montant réservé pour couvrir l'incertitude des estimations (Ch.~\ref{chap:financer}).
    \item[Registre de conformité] Liste des exigences applicables avec leur responsable, leur état et leur pièce (Ch.~\ref{chap:maitriser}).
    \item[Réversibilité] Obligation faite à un prestataire de restituer les données et d'assister la reprise en fin de contrat (Ch.~\ref{chap:financer}).
    \item[Scénario de repli] Portefeuille réduit prévu d'avance pour le cas où une ressource manquerait (Ch.~\ref{chap:financer}).
    \item[Schéma directeur] Document approuvé par l'organe de gouvernance qui fixe pour plusieurs années la cible du système d'information, les projets, leur ordre, leur coût et leur pilotage (Ch.~\ref{chap:vocabulaire}).
    \item[Schéma directeur glissant] Schéma directeur dont l'horizon est repoussé d'un an à chaque révision annuelle (Ch.~\ref{chap:piloter}).
    \item[Stratégie numérique] Choix par lesquels une organisation emploie le numérique pour atteindre ses objectifs ou les redéfinir (Ch.~\ref{chap:vocabulaire}).
    \item[Trajectoire] Ordre daté des projets retenus, lisible année par année (Ch.~\ref{chap:portefeuille}).
    \item[Urbanisation] Démarche qui organise le système d'information en blocs aux frontières stables, à la manière d'un plan de ville (Ch.~\ref{chap:cible}).
    \item[Valeur acquise] Budget des travaux réellement achevés à une date. Elle se compare à la valeur planifiée et au coût réel (Ch.~\ref{chap:piloter}).
    \item[Valeur actuelle nette] Somme des flux d'un projet ramenés à la date de départ, investissement déduit (Ch.~\ref{chap:financer}).
\end{description}
